\subsection{Extending the paired action to every exit arrow}

\begin{result}{E.39}{Proved}{Every GT parameter acts on the full labelled completed exit system}
For every $A\in\widehat{\mathrm{GT}}$, the paired quilt action of
E.34 extends to an equivalence of the strict finite-congruence
completed labelled exit system, with its relabeling, clutching and
forgetful maps.  The equivalence and its operation comparisons
are specified below; changing the auxiliary quilts gives the
specified natural equivalence.  These equivalence classes form an
action of $\widehat{\mathrm{GT}}$, and restriction to genus zero
returns $A$.  Thus restriction from this full-system automorphism
group onto the standard genus-zero GT factor is surjective.

This assertion concerns the ordinary labelled exit system and its
listed operations.  The additional branch-order covering map of
E.37 has the exact compatibility condition stated there.  The
kernel of restriction on all exit arrows is not asserted to be
trivial here; E.35 detects its induced open-group outer actions.
\end{result}
\begin{proof}
Work first with actual multicurves on a fixed connected surface;
transport to a finite orbit skeleton by the explicit functor
(E.145).  For an adjacent incidence $\sigma\supset\tau$, write
\[
 G_0=\Stab(\sigma)/\Delta_\sigma,\qquad
 G_1=\Stab(\tau)/\Delta_\tau,\qquad
 K=\bigl(\Stab(\sigma)\cap\Stab(\tau)\bigr)/\Delta_\tau.
                                                               \tag{E.150}
\]
The distinguished crossing is $e=[1]$.  Its two endpoint maps
are the quotient $p_0:K\to G_0$ and inclusion $p_1:K\to G_1$.
By E.32 its completed incidence orbit is
\[
  (\widehat G_1\times\widehat G_0)/\widehat K,
  \qquad (g_1,g_0)e=g_1 e g_0^{-1},                            \tag{E.151}
\]
where the stabilizer map is $(p_1,p_0)$.  In particular
$\widehat K$ is the full completion of the actual incidence
stabilizer, whether or not either individual endpoint map is
injective after completion.

Choose a complete quilt $Q$ whose pants curves include every
curve of $\sigma$, and pin the branches of the incidence.
Cutting along $\sigma$ gives bordered components with their
restricted quilts.  Their boundary-fixed mapping class groups,
the integer twists on the cut collars, and the finite permutations
of isomorphic components and branches present $K$.  One kills
the twists at the retained nodes $\tau$ and keeps the twist at
the one node to be smoothed.  The source map kills that last
twist as well.  The target map sews its two boundary circles
and then performs the same retained-node quotients.  These are
the actual maps (E.150): extending the cut diffeomorphisms over
the collars gives $p_1$, and normalizing the remaining node gives
$p_0$.  Restriction to the cut pieces is the inverse description;
the common boundary rotations account exactly for the displayed
twist relations.  Branch permutations are retained, with their
actual action on the pieces and collars.

Apply E.34 to these bordered groups and their morphism torsors.
It commutes with every sewing and with branch relabeling, and
sends each boundary twist to its common $\lambda(A)$-power.
Consequently it preserves the twist relations and quotients just
described.  E.36 supplies all marked caps.  We obtain continuous
automorphisms, in the $Q$ trivializations, satisfying
\[
 \Phi_i^Q p_i=p_i\Phi_K^Q\quad(i=0,1).                      \tag{E.152}
\]
This also treats finite component permutations: the functor on
morphism torsors, rather than only its restriction to pure
component groups, is natural for their boundary relabelings.
For a pair of pants the asserted map is the quilt boundary-twist
map; for an annular unit both boundary twists are the same
positive collar twist, as in E.35.  Ordinary marked caps kill
that twist when the collar bounds a marked disc.  Thus these
low-complexity pieces require no additional automorphism choice.

Equation (E.152) defines a map of the entire adjacent biset,
not merely of its pinned stabilizer:
\[
  g_1 e g_0^{-1}\longmapsto
                  \Phi_1^Q(g_1)e\Phi_0^Q(g_0)^{-1}.          \tag{E.153}
\]
The two representatives of a coset differ by $(p_1(k),p_0(k))$;
(E.152) sends that difference to the same diagonal image of
$\Phi_K^Q(k)$.  This proves that (E.153) is well-defined.
It is continuous, and its inverse is supplied by $A^{-1}$.

Fix one endpoint quilt trivialization for each skeletal object.
Let $T_i$ be the quilt comparison from the $Q$ trivialization at
endpoint $i$ to this chosen one, using the orientation of
(E.114)--(E.116), so the endpoint action is
$\Phi_i=\operatorname{Inn}(T_i)\Phi_i^Q$.
In these fixed coordinates the image of the pinned crossing is
\[
                e\longmapsto T_1 e T_0^{-1}.                 \tag{E.154}
\]
Equivariance with $\Phi_1,\Phi_0$ specifies the image of every
other arrow in its incidence orbit.

We verify independence of the refining quilt.  Two complete
quilts refining $\sigma$ are connected by pants moves inside
the cut components and by half-adjustments on the retained
collars.  Indeed the pants theorem applies independently on
each cut component, and the quilt adjustment torsor is the one
proved in E.34.  Every transport weight of such a path lies in
the full cut stabilizer of $\sigma$ and $\tau$.  Its class is
therefore an element $k\in\widehat K$.  Sewing and capping
compatibility identify its two endpoint images with $p_1(k)$
and $p_0(k)$, respectively.  Changing $Q$ therefore changes the
two comparisons in (E.154) by this same diagonal element;
(E.151) makes its image unchanged.  Changing a chosen endpoint
quilt, instead, conjugates the endpoint action and changes all
incident arrows by its one endpoint comparison.  Those
comparisons are precisely the components of a natural
isomorphism.  Their composition is the transport composition
law of E.34, so the specified identifications are coherent.

For a closed ambient surface, choose an auxiliary disc in every
component of the surface cut along $\sigma$.  Let their auxiliary
labels be permuted exactly when those components are permuted,
and use the relabeling morphism torsors of E.34.  Every element
of $K$ lifts to this auxiliary groupoid: represent its cut-component
diffeomorphisms and isotope each to carry the chosen source disc
to the chosen disc of its image component.  Sew with these discs
retained, then fill all of them.  This gives the actual endpoint
maps, including the empty target multicurve.  The filling kernels
are generated by the corresponding boundary twists and disc
pushes; E.36 preserves them and identifies the descended quilt
transports.  The order of filling is immaterial by E.36.  Compare
two systems of discs by retaining their union, after an isotopy
making them disjoint, and filling either subsystem first.  The
same order coherence supplies that comparison.  Thus all component
permutations are retained, without selecting a section of a
capping quotient.  Equations (E.152)--(E.154) descend with these
comparisons; changes of the remaining quilts give exactly the
endpoint gauges already specified.  For a closed torus the
quotient is the explicit $B_3\to\mathrm{SL}_2(\mathbb Z)$
quotient by the square of its center from E.36; genus-zero
unstable pure quotients are trivial.  Hence the construction
also covers the closed and low-complexity endpoint cases without
selecting an unproved section of a capping map.

It remains to check every relation between adjacent arrows.
Take an actual pinned diamond obtained from $\sigma$ by deleting
two distinct curves $c,d$.  Choose one complete quilt refining
this finest multicurve $\sigma$, and restrict or sew it at all
four vertices.  In these simultaneous coordinates, (E.153)
fixes all four pinned arrows; both composites are the same
pinned deletion of $c,d$.  After passage to the fixed endpoint
quilts the comparison at each intermediate vertex cancels:
\[
 (T_2e_{12}T_1^{-1})(T_1e_{01}T_0^{-1})
                 =T_2(e_{12}e_{01})T_0^{-1}.                 \tag{E.155}
\]
The other route has the identical expression.  The pins are
those of this one actual multicurve and are transported by
(E.145), exactly as in (E.147).  Equivariance proves the same
equality for every endpoint translate; continuity proves it
on each completed transport family.  Choice independence just
proved permits any other chosen refining quilts.  Thus all
relations in the complete presentation E.38 are satisfied.
Its universal finite-quotient comparison (E.148)--(E.149)
therefore extends these assignments uniquely to a continuous
functor on the whole completed exit category.  Explicitly, compose
the assignments with each finite categorical quotient of the
target and apply (E.148).  The resulting finite extensions commute
with refinement; their inverse limit is the asserted continuous
functor.

The operation comparisons can be checked in these same
coordinates.  Relabeling relabels the quilt and its pinned
branches.  Clutching takes the quilts on the input components,
inserts the prescribed collar, and uses their sewn quilt;
E.34 gives the identical transport on each component and on the
new collar.  For forgetting a marked point, apply the disc-cap
surgery of E.36 before taking the stabilized graph.  A curve
which becomes peripheral has its twist killed.  Two surviving
collars which become isotopic are stitched, and their integer
and half-integer displacements add.  These are exactly the
maps $q_e=q_iq_mq_j$ and the associated sums of meridians in
(E.123), with the analogous one- and two-collar instances.
The adapted quilt surgery erases the capped displacement and
adds the stitched ones, with zero extra winding, as proved in
E.36.  Accordingly the images of both a pinned incidence and
its two endpoint maps agree with the incidence, identity, or
composite obtained by the actual forgetful functor.  Common
refining quilts check this equality before returning to the
chosen endpoint coordinates; (E.154) supplies its natural
comparison afterwards.  Repeated surgery is concatenation of
the same collar strips, so different orders give the same
comparison.  Multiple sewings use the same simultaneously
sewn quilt.  The sewing/capping comparisons thus respect
composition, disjoint operations and self-clutching.  This
proves compatibility with all the operations stated in the
result, including their specified coherence maps.

Finally, work with the same refining quilt for parameters
$A,A'$.  The group law for the endpoint functors and the
transport cocycle of E.34 imply the group law on (E.153) and
(E.154).  They then imply it on every arrow by E.38.  The
identity parameter gives the identity functor and $A^{-1}$
gives its inverse.  Changing endpoint coordinates yields the
natural equivalences already described, so this is an action
on equivalence classes of full-system equivalences.  On the
open four-point sphere its formula is
$x\mapsto x^\lambda$, $y\mapsto f(x,y)^{-1}y^\lambda f(x,y)$,
with the same actual genus-zero sewing maps as E.34.  Its
restriction is therefore exactly $A$, proving surjectivity.
No assertion concerning a kernel on the adjacent bisets was
used in constructing this right inverse.
\end{proof}
